\section{Formal reading: heads as causal variables}
\label{sec:formal}

Our steering framework has a natural formalization in the language of
causal abstraction~\cite{rubenstein2017,beckers2019,geiger2021,geiger2025}.
A causal abstraction relates two causal models through a mapping between
their states and a mapping between their interventions. Making both
models explicit clarifies what each evaluation metric in our battery does
and does not estimate.

\paragraph{Definitions.}
The high-level model is the Ontologizer's forward pass, treated as a
causal model over the input and the heads: each head $(\ell,i)$ is a
candidate \emph{variable} with value space $\{1,\dots,k\}$, the input
$\mathbf{x}$ is a parent of every head, and under residual forwarding
every earlier head is a parent of every later head because layer
$\ell{+}1$'s classifier takes as input the residual produced by layers
$\le \ell$. These causal edges are internal to the code: unlike
transcoder circuits~\cite{dunefsky2024}, which pass through a base
model's own sublayers, they describe how the code iteratively refines a
single representation rather than how the pretrained encoder computes it.
Set-interventions applied via \texttt{withArgs} serve as hard
interventions in this model: $\mathrm{do}[(\ell,i){=}j']$ forces the
head's assignment while leaving all non-descendants unchanged by
construction, and marginalizing out the continuous intermediate states of
the forward pass defines an exact transformation by
construction~\cite[Thm.~10]{rubenstein2017}. Because a soft head's state
is a point on its simplex, discretizing it into one of $k$ categorical
values constitutes a lossy value merge~\cite{geiger2025}, so forcing an
entry on a soft head acts as a steering intervention rather than an
interchange intervention. The low-level model is the pretrained encoder
that produces the input embeddings, and the mapping between states is
$\mathrm{assign}$, which maps an embedding $\mathbf{x}$ to its per-head
assignments. Whether the head model abstracts the encoder under
$\mathrm{assign}$ is an approximate-abstraction
question~\cite[Def.~41]{geiger2025}, and the abstraction cannot be
constructive~\cite{beckers2019} because every head takes the same
embedding as input.

\paragraph{What the instruments estimate.}
Input-space perturbations are not interventions that such an abstraction
maps; instead, they measure controllability of the head model rather than
whether the encoder computes using those
heads~\cite[Sec.~3.7]{geiger2025}. Specifically, the realization rate in
\texttt{steerembed.py} tests a single coordinate of the condition that a
perturbed input's code matches the code under
$\mathrm{do}[(\ell,i){=}j']$, so any realization rate below $1$ rules out
an exact transformation mapping that perturbation to the target
intervention. Its collateral metric counts the remaining heads whose
assignments change relative to the \emph{unsteered} code, which in a
residual stack includes the downstream reclassification that
$\mathrm{do}[(\ell,i){=}j']$ itself induces. Similarly, the
cycle-consistency gain in \texttt{steerfid.py} evaluates a text-level
implementation across the closed loop $\mathrm{assign} \circ
\mathrm{encode} \circ \mathrm{decode}$. Because it patches nothing from a
separate source input, it measures a normalized activation gain net of a
random control rather than interchange-intervention
accuracy~\cite{geiger2022}, and the loop's ${\sim}13\%$ unsteered
round-trip survival rate (Section~\ref{sec:steer}) caps how often an
exact intervention could be detected through text. Cross-seed matching
tests whether the $\mathrm{assign}$ maps of two independent training runs
agree up to entry permutation. In principle, the softmax parameterization
leaves per-head entry ordering and any logit shift or scale absorbed by
the temperature unidentified; consequently, for heads to reflect
intrinsic structure of the data rather than artifacts of a single run,
both the sample partitions and the decoded directions must reproduce up
to entry permutation. Section~\ref{sec:seeds} tested both criteria on
hard codes across SONAR and GPT-2. Beyond layer~0, no individual head
reproduces, nor do our layer-level diagnostics show cross-seed agreement
beyond each model's shared alignment with the data's principal
directions. Testing whether a full layer reproduces as a variable merge of its
heads~\cite[Def.~37]{geiger2025}, in which individual heads form
seed-dependent factorizations, requires a whole-layer comparison that
translates one seed's layer code into another's. Layer~0 reproduces only
partially, and on SONAR almost entirely through a single large-atom head.

\paragraph{Why the cascade costs reproducibility.}
Under this causal framing, the residual stack's lower cross-seed
reproducibility at matched head size (Section~\ref{sec:l0size}) and its
higher input-space intervention collateral (Section~\ref{sec:entangle})
arise from the same mechanism. Every head computes a function of its own
layer input, and every head past layer~0 takes as input the residual left
by upstream heads; consequently, any change in the upstream decomposition
alters the input distribution that downstream heads partition. Even
layer~0, which operates directly on the raw input, is only partially
constrained: because its heads are trained jointly with downstream
layers, which partition of the input minimizes total loss depends on how
subsequent layers model the remaining residual. A single-layer
architecture avoids both dependencies: at matched head size, its heads
reproduce more consistently across seeds and realize input-space
interventions far more reliably, with half the collateral. Nevertheless,
single-layer hard heads are still not independently addressable, because
hard-code classifiers lie close to their decision boundaries in both
stacked and flat architectures; residual depth therefore compounds an
entanglement already present within each layer.

\paragraph{Metric hygiene.}
Geometric claims about decoded directions should not depend on an
arbitrary choice of inner product, a point emphasized by Park et al.\ in
their analysis of language-model representations under an estimated
causal inner product~\cite{park2025}. Our within-head repulsion finding
is invariant to this choice: under the inverse-variance inner product of
the whitened training objective (which whitens the data rather than
estimating a causal inner product), within-head repulsion in the softmax
Ontologizer strengthens ($\zeta_{\cos}$ of $-24.4$ vs.\ $-17.6$ under the
Euclidean inner product), while the coherence profiles of both comparison
models remain unchanged. However, this measurement does depend on the
parameterization: because a soft head's assignments sum to one, its
decoded directions are determined by the forward function only up to a
shared per-head offset, meaning that uncentered cosines characterize the
trained atom parameters rather than the input--output function alone
(Section~\ref{sec:prelim-orth}). Two additional methodological rules
emerged from our analysis. First, geometry must be evaluated in output
space rather than in dictionary space, whose decoder null-space component
is unconstrained by reconstruction (Section~\ref{sec:rowcos}) and
represents the same kind of causally disconnected subspace that can
render subspace patching misleading~\cite{makelov2023}. Second, a layer's
scalar residual gain must be separated from its atom norms before
attributing effects to head magnitude: across layers, differences in
output contribution magnitude reflect the layer-wide residual gain rather
than larger individual atoms (Section~\ref{sec:l0size}).

\paragraph{Composition.}
Interpreting heads as causal variables provides little benefit if
interventions on two heads interact unpredictably. Using
\texttt{compose.py}, we measure the additivity of decoded intervention
deltas on \texttt{resid\_nc} under the whitened metric, reporting medians
across head pairs:

\begin{center}
\begin{tabular}{@{}lrrrr@{}}
\toprule
stratum & pairs & cos & rel.\ err & later layers \\
\midrule
same-layer L0 & 28 & 0.959 & 0.350 & 4 \\
same-layer L1 & 28 & 0.947 & 0.354 & 3 \\
same-layer L2 & 28 & 0.967 & 0.259 & 2 \\
same-layer L3 & 28 & 0.996 & 0.092 & 1 \\
same-layer L4 & 28 & 1.000 & 0.019 & 0 (numerical floor) \\
cross-layer (earlier $=$ L0/L1) & 42 & 0.97--0.98 & 0.20--0.23 & \\
cross-layer (earlier $=$ L2/L3) & 18 & 1.00 & 0.02--0.04 & \\
\bottomrule
\end{tabular}
\end{center}

Three patterns emerge from these measurements. (i)~Joint two-head
interventions align closely with the sum of the corresponding single-head
interventions across all strata (median cosine $\ge 0.95$), confirming
that \texttt{withArgs} set-interventions compose approximately
additively. (ii)~Deviations from exact additivity arise entirely from
downstream reclassification: relative error decreases monotonically with
the number of remaining downstream layers and reaches the numerical floor
in the final layer (where the measured $0.019$ reflects float32 rounding
amplified by decoder cancellation, dropping to ${\sim}3{\times}10^{-6}$
in float64). (iii)~Cross-layer pairs compose \emph{more} additively than
same-layer pairs at matched depth because a set-intervention overrides
upstream perturbations on the later head, making cross-layer interaction
strictly one-way. This behavior contrasts sharply with input-space
perturbations (Section~\ref{sec:entangle}): \texttt{withArgs} sets a
head's assignment directly and composes cleanly, whereas an
\emph{input-space} step perturbs every head simultaneously and flips most
of them. Downstream reclassification, which introduces only a 5--35\%
deviation under direct code-level interventions, becomes the dominant
effect of an input-space step.

\section{Preliminary conclusions}

\begin{enumerate}
\item \textbf{The trained atoms form contrast sets.} In the softmax
Ontologizer, within-head decoded directions repel relative to random
subsets in 96\% of heads, and in the hard-code stack, within-head decoded
atoms are roughly orthogonal to slightly negatively correlated, so each
head's entries represent well-separated alternatives. Because the forward
function determines within-head atoms only up to a shared per-head
offset, this repulsion characterizes the trained parameters; isolating
the functional contribution requires comparing atoms centered within each
head against a similarly centered null. Neither reconstruction metrics
nor per-latent probes capture this contrast-set structure.
\item \textbf{Between-head coordinate partitioning does not emerge in
summed models.} No trained summed model partitions its dictionary
coordinates across heads. Between-head support overlap sits at or above a
coordinate-shuffle null across all summed runs, each head's usage spans
52--98\% of dictionary coordinates, and this overlap persists in the
subspace projected through the decoder as well as in its null space. Only
very long training (8.27M steps) reduced overlap, doing so by
sparsifying individual head profiles rather than separating heads onto
disjoint coordinates. Consequently, static weight-level analysis cannot
assume disjoint heads in summed models; concatenated heads remain the
only mechanism that guarantees disjoint support, doing so by
construction.
\item \textbf{A discrete ontology is viable, though not by discretizing a
softmax model.} Taking the one-hot argmax of a trained softmax
Ontologizer fails to reconstruct the input, yielding at least $77\times$
higher error than the 800-bit Gaussian reference on \texttt{resid\_nc}
and nine orders of magnitude higher error on \texttt{softmax\_shm}. By
contrast, training a straight-through arm with $2.4\times$ as many heads
makes the discrete code the forward pass itself: with its four per-layer
gains pinned to their held-out means, it forms a purely discrete 1900-bit
code at $\fvu = 0.158$. At a matched 1900-bit rate, neither the converged
$k_{\mathrm{SAE}}{=}160$ $\topk$ SAE nor the single-layer hard variant
can improve on $\fvu = 0.236$--$0.273$ under any coefficient
quantization. However, a coarsely quantized, entropy-coded L1 SAE reaches
$\fvu = 0.129$ at an average rate of 1900 bits per row and $0.152$
(comparable to the hard code) when capped at 1900 bits on every row;
averaged across rows, the hard code outperforms all SAEs only under
fixed-length support indexing ($0.158$ vs.\ $0.180$), because its entries
are used nearly uniformly whereas an L1 SAE's latents have skewed firing
rates. On the sentence-split cache, a wider $\topk$ SAE
($k_{\mathrm{SAE}}{=}256$) outperforms the hard stacks under both average
and per-row bit budgets ($0.128$ and $0.129$ vs.\ $0.161$ and $0.164$;
Appendix~\ref{sec:quantrate}). The hard code sits $3.1\times$ above the
Gaussian reference for its rate, and so at least that far from the data's
distortion--rate function, a limit over long blocks of samples that a
code describing each sample alone need not be able to
reach~\cite{cover2006}. Training this discrete stack required matching
the classifier temperature to the logit scale, applying relative logit
noise, rescaling the weight on the batch-mean uniformity term
$\mathrm{KL}_m$, and rescaling dictionary initialization (which alone
reduced held-out error by one third).
\item \textbf{The discrete code preserves downstream model performance.}
On GPT-2 small, an 800-bit hard code with signed direct atoms in
activation space and per-layer gain--shape normalization recovers 97.6\%
of next-token cross-entropy loss, surpassing a $k_{\mathrm{SAE}}{=}32$
SAE (96.2\%) at nearly the same unwhitened $\mathrm{FVU}$, while
increasing head width to 64 entries surpasses that SAE on both metrics.
Across model families, $\mathrm{FVU}$ does not monotonically order
downstream loss recovery. Training continuous low-rank fibers on top of a
frozen discrete base recovers nearly all remaining loss (99.9\%) while
leaving the discrete base assignments unchanged, showing that a discrete
ontology and near-lossless reconstruction can coexist at the cost of
hundreds of additional continuous coordinates per token.
\item \textbf{For soft codes, reconstruction measures structural cost
rather than steering utility.} Unstructured codes achieve roughly
$50\times$ lower reconstruction error at matched coefficient budgets,
with the per-head simplex constraint alone accounting for roughly a
$30\times$ penalty. Downstream text fidelity on SONAR is a monotone
function of $\fvu$ and adds no independent evaluation axis.
\item \textbf{Residual depth improves reconstruction at the cost of
addressability and reproducibility.} Compared to a single-layer model at
equal bit budget, parameter count, and near-uniform entry usage, the
five-layer stack achieves $1.79\times$ lower reconstruction error, driven
by conditioning each stage on the sample's remaining residual error
(rather than by utilization, input conditioning, or the scalar gain
side-channel, though deep supervision on prefixes has not yet been
ablated separately). Conversely, an input-space perturbation flips 66\%
of unrelated heads in the five-layer stack vs.\ 33\% in the single-layer
model, and realizes its target entry on only 37\% of rows vs.\ 99.6\%; in
both models, off-target collateral matches that of a norm-matched random
step because hard-code classifiers sit close to their decision
boundaries. Moreover, at matched head contribution size, stacked heads
show several-fold lower cross-seed reproducibility than single-layer
heads across every layer, including layer~0. Although the single-layer
softmax variant \texttt{g160softmax} matches the reference checkpoint
\texttt{resid\_nc} in reconstruction, a configuration-matched softmax
Ontologizer achieves $6\times$ lower error both on out-of-sample rows
that no model saw during training and on the evaluation tail, indicating
that \texttt{resid\_nc}'s lack of a depth advantage was
checkpoint-specific.
\item \textbf{Steering succeeds in embedding space, whereas text
round-tripping obscures it.} Stepping along a decoded atom at strength
$0.25$ realizes the target entry on 82\% of rows in the softmax model
(vs.\ a 3\% chance rate), whereas decoding and re-encoding through text
preserves unsteered entries on at most ${\sim}13\%$ of inputs, bounding
how often exact interventions remain observable after round-tripping. At
low collateral, the soft code is more steerable than the hard code.
Steerability does not track reconstruction error, and the classifier's
input gradient provides an effective steering direction only in layer~0,
where the raw input enters the classifier directly.
\item \textbf{Individual heads beyond layer~0 do not reproduce across
seeds, and layer~0 reproduces only partially.} Independent seeds converge
to identical held-out reconstruction error to four decimal places while
learning substantially different reconstructions (whitened residual
correlation of $0.30$, compared to $0.83$ between SAE seeds). Past
layer~0, independent runs show no above-null head-level agreement across
partitions, output contributions, subspaces beyond the data's principal
directions, or dictionary Gram matrices on either SONAR or GPT-2
(including under the direct-atom design of Section~\ref{sec:gpt2hard}),
even though individual heads remain stable across late checkpoints within
a single run. Layer~0 reproduces at less than one fifth of the within-run
control. Its higher cross-seed agreement relative to deeper layers does
not stem from larger atom norms: its atoms are no larger than those of
later layers, and its larger output contributions reflect the layer-wide
residual gain. Within layer~0, however, heads with unusually large atom
norms account for the reproducible signal---on SONAR, almost entirely a
single head. This lack of identifiability is the central limitation for
our steering framework: a head that does not recur across runs cannot be
assigned a stable meaning, qualifying every per-head claim in this study.
The sole exception to date is a block of heads explicitly trained for
temporal consistency: on pairs of sentences from one document, a Fano mutual-information
objective causes a block of 38 slow heads to reproduce across three seeds
(Hungarian-matched NMI of $0.412$--$0.414$ vs.\ a $0.005$ null) while
varying more slowly across sentences than any linear projection of the
embedding. Even there, reproducibility holds at the block level rather
than for individual heads: two slow heads within the same seed agree at
NMI $0.40$, so individual slow heads are not yet uniquely identified
across runs (Section~\ref{sec:slowheads}).
\item \textbf{Entries contribute nearly equally except when classifying
the raw input.} In layers~1--3, entries within each head carry nearly
equal shares of output contribution (Gini $0.06$--$0.15$, with zero dead
entries), whereas layer~0's entries carry unequal shares (Gini
$0.33$--$0.59$) because its atom norms vary widely within heads, despite
near-uniform entry usage. Because the
single-layer hard model exhibits the same dispersion, this inequality
arises from classifying the unconditioned raw input rather than from
residual stacking.
\item \textbf{Learned heads prioritize semantic content over source
language.} Consistent with SONAR's design as a language-agnostic semantic
embedding, the strongest head (by $\eta^2$) partitions inputs by broad topic and
genre, and is the single head that recurs across random seeds. Source language, the corpus's primary free ground-truth label, served as a
diagnostic probe rather than an expected single-head variable: language
identity is distributed across multi-head neighborhoods rather than
isolated in a single head, reflecting the weak separation of languages in
SONAR space. On that probe, the hard code captures $3.5$--$3.9\times$ as
much language information as the softmax reference, confirming that
learned partitions align with known attributes above chance; directly
testing whether heads recover natural semantic variables will require
topic and genre annotations.
\item \textbf{Per-latent metrics mischaracterize classifier codes.}
Standard sparse probing ranks the Ontologizer lowest even where
partition-level analysis identifies verified causal language structure.
Across models, per-feature autointerp detection F1 is roughly 80\%
explained by activation density, and even more strongly by how many
languages a feature's top activating examples span, which covaries with
density. When restricted to features spanning as many languages as random
draws, the sparse SAE's F1 lead over the Ontologizer narrows to $0.10$ or
less. Furthermore, parameter-only decodes generate single example
sentences rather than category descriptions and become legible only after
subtracting the head's mean atom, consistent with simplex models of
categorical concepts in which value vectors share a parent offset and
separate only through their mutual differences~\cite{park2025}.
\item \textbf{Non-negativity substantially degrades dictionary geometry
and reconstruction.} Constraining dictionary weights to the non-negative
orthant induces high atom collinearity, reduces effective rank to one
third of a signed dictionary's, and increases held-out $\fvu$ by 26\% at
the wide dictionary and by 90\% with a square decoder; mean within-head
atom cosine exactly orders all five GPT-2 orthant arms by held-out error
at convergence. Signed atoms eliminate this collinearity and achieve
lower error at both dictionary widths. Moreover, dictionary collinearity
statistics computed prior to decoding largely reflect unconstrained
decoder null-space components.
\item \textbf{Every diagnostic in this architecture requires a
shape-matched null.} Non-negativity constraints and layer-dependent
residual gains induce nontrivial baseline values across geometric and
information-theoretic metrics. Three of our four initial cross-seed
reproducibility metrics, the apparent HSIC between-head dependence, raw
dictionary-space collinearity measurements, and the apparent effect of
head magnitude on cross-seed reproducibility all reversed or vanished
once compared against shape-matched nulls.
\end{enumerate}

\section{Follow-up investigations}
\label{sec:followup}

We order the primary open investigations by their importance to
unsupervised steering and auditing.

\begin{itemize}
\item \textbf{Head identifiability first.} Before individual heads can be
interpreted or used in audits, we must either make them reproducible
across training runs or identify the subset of heads that reproduce.
Cross-seed partition comparisons on the softmax Ontologizer at 370k steps
show virtually no head-level agreement (Appendix~\ref{sec:softseed}). A
controlled $k{=}128$ experiment sharpens the source of this sensitivity:
two distinct architectures trained from the same random seed (one with a
learned decoder and one with direct atoms) learn matching partitions head
for head (matched NMI $0.18$, with every head pairing with the co-indexed
head), whereas each architecture compared across two seeds reaches only
$0.064$ vs.\ a $0.059$ null. Because seed replicas iterate through the
same training examples in the exact same order, the random seed acts
solely through parameter initialization and training noise; those two
factors, rather than the dataset or downstream parameterization past the
classifier, largely determine how each head partitions the space, echoing
findings in toy networks where identical training data yield equivalent
algorithms implemented via distinct features~\cite{chughtai2023}.
Classical identifiability theory points to two remedies. The first is
unsupervised model and feature selection via cross-seed agreement, as in
Unsupervised Disentanglement Ranking, which scores models by median
cross-seed alignment without ground-truth labels~\cite{udr2020} (in
latent-variable models, three or more conditionally independent views of
a discrete variable identify its distribution up to permutation when the
class-conditional distributions satisfy a rank
condition~\cite{allman2009}, though here individual heads serve as
candidate learned views rather than fixed ones, and standard sparse-coding
uniqueness guarantees require lossless
reconstruction~\cite{hillar2015}). The second is breaking static marginal
symmetry using paired observations, where temporal
contrast~\cite{pcl2017}, sparse transition priors~\cite{slowvae2021},
shared factors across corruptions~\cite{locatello2020}, auxiliary
conditioning~\cite{ivae2020}, or contrastive conditional
alignment~\cite{zimmermann2021} guarantee identifiability up to
permutation under their respective generative assumptions. Initial
experiments on the sentence-split SONAR cache \texttt{mc4\_sent16}
(Appendix~\ref{sec:seedsdetail}) evaluate this second approach. Linear
ridge regression from $\mathbf{x}_t$ to the next sentence
$\mathbf{x}_{t+1}$ reaches $\mathrm{FVU}_{\mathrm{next}}{=}0.896$
(roughly 10\% of next-sentence variance explained linearly, which does
not upper-bound nonlinear predictors), and adjacent sentences exhibit
strong directional alignment (raw cosine AUC $0.923$ vs.\ $0.685$ in
whitened distance, where norm variance dominates), yet heads trained on
single sentences switch entries on 92.2--98.1\% of adjacent pairs and
retain their entry barely above chance (Cohen's kappa at most $0.081$
across all heads). Directly penalizing cross-sentence entry switches on
38 designated slow heads collapses their usage entropy to 0.04--0.09 bits
for penalty weights $s_{\mathrm{pair}} \ge 10^{-4}$. Replacing that
penalty with a Fano lower bound on each slow head's cross-sentence mutual
information prevents collapse and yields a 38-head slow block that
reproduces across three seeds, though not yet head by head
(Section~\ref{sec:slowheads}). Immediate next steps for paired-view training include decorrelating the
slow block (whose heads currently concentrate along the top one or two
linear slow directions and agree with one another nearly as strongly as
they agree across seeds, preventing single-head identification),
extending the objective to triples of spaced sentences within a document
to supply the three conditionally independent views required by classical
rank conditions~\cite{allman2009}, characterizing what the slow heads
encode beyond source language, and applying the paired-view objective to a
causal language model over multi-turn conversation states. Parallel
single-sentence reproducibility checks remain open as well: training
full-length (4.4M-step) cross-seed replicas of \texttt{resid\_nc\_hm} and
of \texttt{g160softmax}; testing whether regularization can extend the
reproducibility observed in large-atom heads; testing whether the binary
($k{=}2$) arm's one-bit heads reproduce more reliably than 32-way heads;
isolating why the layer-0 topic head recurs when other heads do not (at
$k{=}128$, a dominant layer-0 head recurs at one third to one half of the
$k{=}32$ topic head's agreement); testing whether per-head router gains
drive agreement, since on GPT-2 adding a router doubles contribution
agreement while deactivating the final two layers; and training a
single-layer GPT-2 arm to replicate the size-matched comparison of
Section~\ref{sec:l0size} on a second substrate.
\item \textbf{Consolidating GPT-2 hard-code features on the main branch.}
Direct atoms have been ported to the main branch
(Appendix~\ref{sec:gpt2ablation}), whereas \texttt{resid\_first} and
\texttt{gain\_clip} have not, so the GPT-2 runs of
Section~\ref{sec:gpt2hard} currently reproduce only on the
\texttt{headline} branch. On SONAR, a $5\times54$ stack with $k{=}128$
entries per head (1{,}890 bits, matching the 1{,}900-bit $k{=}32$ stack)
converges to roughly twice the error ($\fvu = 0.31$--$0.32$ vs.\
$0.1535$, two seeds per arm), with direct atoms outperforming decoded
atoms by $0.009$. Because both $k{=}128$ arms achieve full entry
utilization, this gap does not reflect dead entries; however, the
temperature, learning rate, and weight on the batch-mean uniformity term
$\mathrm{KL}_m$ were carried over directly from $k{=}32$ and must be
retuned before attributing the degradation to $k$. Direct atoms also make
set-interventions act directly in activation space by adding
$g_\ell(\mathbf{a}_{ij'} - \mathbf{a}_{ij})$ at layer $\ell$, an
intervention whose steering behavior remains to be measured.
\item \textbf{Per-layer decoders vs.\ width-matched shared decoders.}
Early one-epoch direct-atom pilots on GPT-2 showed less than a $0.007$
improvement in unwhitened $\mathrm{FVU}$ from separate per-layer
decoders, initially suggesting little benefit. On a twin of the 40M-token
concatenated-head model (Appendix~\ref{sec:gpt2ablation}), however,
giving each layer its own decoder reduces held-out $\mathrm{FVU}$ from
$0.186$ to $0.169$ ($97.3\% \to 97.7\%$ loss recovered) and reduces
cross-layer coupling between co-indexed heads from $97\%$ to $6\%$ (near
the $3\%$ chance floor), demonstrating that the strong coupling arose
from projecting head $i$'s slice through the same shared decoder block
across all five layers. Because using $l{=}5$ separate decoders
multiplies decoder parameter count fivefold, a shared-decoder control
matched in total decoder width ($d_{\mathrm{head}}{=}240$) is still
required.
\item \textbf{Two-stage fiber training on SONAR.} Training low-rank
fibers on top of a frozen discrete base preserved the base code's
assignments and error on GPT-2 while matching the overall reconstruction
of unconstrained joint fiber training. Although fibers are implemented on
the main branch, no SONAR run has yet used them; this two-stage recipe
should be evaluated on SONAR against SAEs at matched continuous capacity
(where rank-1 radial fibers require 160 continuous values per token).
\item \textbf{Re-evaluating ablations under the rescaled initialization.}
Our input-space entanglement diagnostics and the gain and blend ablations
were run under the original dictionary initialization, which started
$1.84\times10^8$ off at step~0 and left 10\% of code
capacity unused; these ablations should be repeated under the rescaled
initialization.
\item \textbf{Scaling the steering evaluation.} We have evaluated
supervised steering vectors against model-derived directions on the
models' own partitions (Section~\ref{sec:steer}) and, for
\texttt{resid\_nc\_hm}, on external language steering
(Appendix~\ref{sec:steeroverlay}). Several steering comparisons remain
open: running the external language-steering benchmark on the hard-code
stack and the single-layer hard arm; testing bilinear eigenfeatures as
steering directions; measuring across-context consistency of entry-swap
effects, which for SAE features is partially predictable prior to
intervening~\cite{duan2026}; and weighing the single-layer hard arm's
near-perfect target realization, halved collateral, and higher cross-seed
reproducibility against its $1.79\times$ higher reconstruction error.
Reducing input-space collateral appears to need larger classifier margins
more than better steering directions. Even with multi-layer
residual reclassification excluded, input-space collateral is largely
direction-blind: linearized within a single layer, stepping along any
Ontologizer's decoded atom flips other heads at the same rate as a random
step of equal norm, and in decoded text, ablating a head's contribution
perturbs outputs no more than a norm-matched random step. Because hard-code classifiers lie close to their decision boundaries (a
within-layer linearized random step of $0.25\lVert\mathbf{x}\rVert$ flips
30\% of same-layer heads), distance from the boundary looks like the
lever. Consistent with this hypothesis, \texttt{ste\_k2\_h380}
(five layers of 380 binary heads) exhibits substantially lower
collateral: the same $0.25\lVert\mathbf{x}\rVert$ step is predicted to
flip 13\% of other heads to first order and empirically flips 17--19\%
across the full stack for all steering directions (random included), vs.\
66\% in the 32-entry stack. Although its raw logit margins are $34\times$
larger than the 32-entry stack's, raw margin scale alone is not the
cause, as raw margins doubled between steps 370k and 739k while the
first-order flip rate remained unchanged because overall logit norms
scaled proportionally. Neither training duration nor the $\mathrm{KL}_m$
weight explains the reduction: at step 370k (matching the 32-entry
stack's full training horizon), both the first-order flip prediction and
empirical collateral were already at 13\% and 17--19\%, and the
batch-mean uniformity weight matches the 32-entry stack's effective
per-head weight (since $\mathrm{KL}_m$ averages across heads, a weight of
$5\times10^{-4}$ over 380 heads applies the same per-head penalty as
$10^{-4}$ over 76 heads). However, because \texttt{ste\_k2\_h380} was
trained on a single seed at half the batch size (and thus under noisier
per-step gradients), a batch-matched multi-seed comparison is still
needed to isolate the effect of $k$.
\item \textbf{Task reconstruction.} A central goal for reproducible heads
is identifying the classification task each head performs. Head-level autointerp, which summarizes a head from side-by-side examples
across all of its entries and scores blinded assignment accuracy against
sibling entries, found no above-chance signal across 20 unscreened heads
per condition (Appendix~\ref{sec:headautointerp}). By contrast, when
scored by graph modularity on a nearest-neighbor graph of each layer's
input residual, partitions in layers~1--4 of the softmax Ontologizer are
two to three times as modular as size-matched random regroupings of their
entries, whereas neither post hoc discovered groups nor the trained heads
of a Matryoshka SAE~\cite{bussmann2025} exceed null modularity on matched
prefix-residual graphs (evaluated on one checkpoint per model;
Appendix~\ref{sec:modularity}). Remaining task-reconstruction experiments
include measuring mutual information between head partitions and metadata
label libraries (prioritizing topic and genre from labeled corpora or
unsupervised topic clusters to match SONAR's primary semantic axes,
alongside length, formatting or noise, and script as secondary
attributes), re-running head-level autointerp on heads that pass
cross-seed reproducibility screening with sufficient sample size to
resolve above-null accuracy, and evaluating invariance across parallel
translations in FLORES to separate semantic-content heads from
surface-form heads. Each protocol should be run identically on stacked
Ontologizer heads and on the single-layer \texttt{g160softmax} heads.
\item \textbf{Hierarchy as MoE routing.} Measuring conditional mutual
information between successive layers' partitions would test whether
residual depth constructs compositional tasks, while evaluating
\texttt{forward="labels"} under the rescaled initialization would test
our residual-conditioning explanation from the opposite direction.
Pairwise effective information~\cite{hoel2013} upper-bounds direct edge
weights in the causal graph of Section~\ref{sec:formal}: in the hard-code
stack, forcing an upstream head to a live entry shifts any single
downstream head's distribution by roughly $0.01$ out of a possible $5$
bits across all depths, though this causal effect is $11$--$16\times$
larger than the pair's observational mutual information. Whether upstream
interventions exert diffuse weak effects or strong joint effects across
multiple heads, which pairwise statistics cannot detect, remains open.
\item \textbf{Mechanism of the collinearity setpoint.} The collinearity
setpoint controller, which improved soft-model steerability, currently
acts on a dictionary-space cosine largely dominated by decoder null-space
freedom; this controller should be tested on decoded atom cosines and on
signed dictionaries.
\item \textbf{Penalizing head overlap in output space.} When trained as a
dictionary-space penalty, support overlap $\omega$ reaches its target
solely by shifting head supports into the decoder's null space, or into
near-null singular directions when the decoder is square
(Appendix~\ref{sec:support}). To take effect on reconstructions, an
overlap penalty must act on decoded atoms or on direct atoms that have no
decoder. Our $k{=}128$ direct-atom runs used $s_\omega{=}10^{-6}$, which
proved too weak to reduce $\omega$ (on GPT-2, $\omega$ responds only from
$s_\omega \ge 10^{-5}$). Because raw activation coordinates are not
privileged axes (product quantizers perform substantially better after a
learned orthogonal rotation than on raw coordinates~\cite{ge2014}),
forcing heads onto disjoint output coordinates is a restrictive and
basis-dependent constraint, though one that cannot be satisfied via
null-space directions. Concatenated heads remain the only reliable route
to disjoint support.
\item \textbf{Logit-scale control for straight-through training.} Without
explicit per-head logit normalization (\texttt{logit\_norm}), classifier
logit norms rescale during training to offset temperature annealing after
an initial transient (Appendix~\ref{sec:gpt2ablation}). Future
straight-through runs should normalize logits per head, regulate
temperature $t$ via a logit-spread$/t$ setpoint controller rather than a
hand-tuned schedule, or add a logit-scale regularizer analogous to the
commitment loss that anchors encoder outputs to codebook vectors in
VQ-VAE~\cite{vandenoord2017}.
\item \textbf{Rate-aware training.} Our operational rate--distortion
comparisons evaluate models that were never trained for entropy coding:
the hard stack uses entries nearly uniformly (maximizing per-head index
entropy), and none of the SAEs were trained with quantized coefficients.
Training both families directly against entropy-coded rate (for example,
penalizing entry-usage entropy in hard codes and sweeping
quantization-aware SAEs across the 1{,}900-bit regime) would compare the
architectures themselves rather than fixed-rate against unquantized
models.
\item \textbf{Preventing dead entries in single-layer hard models.} The
29\% dead-latent rate in \texttt{g160top1} stems from within-head
winner-take-all starvation, which persisted despite enabling Gao et al.'s
dead-latent auxiliary loss~\cite{gao2024} (\texttt{aux\_k} of 512
throughout training), because auxiliary reconstruction gradients do not
force a starved latent to win its head's argmax. Alongside winner
dropout, two established remedies warrant testing: annealing per-head
selection from top-$k_{\mathrm{sel}}$ down to top-1 over the first half
of training, analogous to k-sparse autoencoders reducing their active
budget from 100 to the target sparsity over the first half of training to
prevent unit starvation at low $k_{\mathrm{SAE}}$~\cite{makhzani2014}, and
re-initializing underutilized entries from current-batch residuals as in
neural audio codecs~\cite{zeghidour2021} (already supported by neuron
resampling in \texttt{sae.py}).
\item \textbf{Training for progressive truncation.} Adding a
deviation-sparsity penalty during soft-model training would test whether
the soft Ontologizer's Pareto frontier can shift toward lower capacity,
analogous to multi-sparsity training for progressive $\topk$ SAE
codes~\cite{gao2024}.
\item \textbf{Completing the structural ablation ladder and baselines.}
Combining head grouping with nested prefixes (\texttt{groups+prefixes},
rungs 2a+3) has now been trained as \texttt{m5120\_g160top1\_p5}
($\fvu = 0.294$) for the modularity control above, though it has not yet
been added to the full evaluation ladder. Remaining baseline runs include
the single-layer variants at $m{=}11264$, a mid-range $\topk$ sweep
($k_{\mathrm{SAE}} \in \{512, 2048\}$), and random-support and
non-negative least-squares refit controls, the latter matching the
positivity-constrained refit of Gao et al.~\cite{gao2024}.
\item \textbf{Closed lines of inquiry.} Several candidate modifications
require no further experiments: enforcing between-head independence via an
HSIC penalty (unregularized heads already partition nearly independently,
and the biased HSIC estimator induces genuine head freezing because only
a constant head escapes its positive bias floor;
Appendix~\ref{sec:freeze}), decoder sparsification and decoder bias terms
(which produce no interpretable structure and alter error by $<1\%$),
forwarding classification codes alongside the residual, explicit input
whitening, and rescaling batch-mean uniformity weights across
architectural shapes (the latter two being subsumed by rescaled
initialization). Finally, the layer-4 head freezing in \texttt{resid\_nc}
is explained by residual saturation and is unchanged by the
evaluation-probe correction, leaving only two anomalously coherent heads
(L2h25 and L3h28) for individual inspection.
\end{itemize}
